Die Projektdauer bedingte die Einteilung des Aufbaus in Maturity Level. Level 1 stellt den Prototypen dar, der innerhalb der Projektzeit realisiert wurde. Dieser Level wurde nahezu vollständig wie geplant abgedeckt von unserem Prototyp. Dies bedeutet, dass die Projektplanung von uns realistisch durchgeführt und festgehalten wurde.

Eine der Änderungen, die am Prototyp vorgenommen werden sollten, betreffen die Nutzung der Aktoren und Sensoren für die Fortbewegung und Orientierung des Roboters. Manche Punkte haben sich nach den Tests als ausreichend bis mangelhaft für eine akkurate Bewegung über den Tisch erwiesen.

Andererseits heißt dies, dass wir bereits in der Prototypenphase erfolgreich Unzulänglichkeiten im Design erkennen konnten, die in den folgenden Maturity Leveln  korrigiert werden können.

Deshalb halte ich die Prototypenphase und somit auch das Projekt im gesamten als erfolgreich abgeschlossen.
